% test-005.tex -- tagged-PDF test de las utilidades de algebra y estadistica
%
% Compilar:  lualatex-dev test-005.tex
% Validar:   show-pdf-tags --xml test-005.pdf | rnv-wrapp
%            verapdf --flavour ua2 --format text test-005.pdf
%
% Test de prueba para las utilidades de álgebra y estadística: las formas de
% cada comando y sus llaves. Los "doc:" son lecturas reales de NVDA 2026.2 con MathCAT.
\DocumentMetadata
  {
    lang=es-CL,
    pdfversion=2.0,
    pdfstandard=ua-2,
    tagging=on,
    tagging-setup={math/setup={mathml-SE}},
  }
\documentclass{article}
\usepackage[spanish]{babel}
\usepackage{unicode-math}
\usepackage{spintent}
\usepackage{hyperref}
\hypersetup
  {
    pdfdisplaydoctitle = true,
    pdftitle           = {test-005: utilidades de álgebra y estadística},
  }

\begin{document}

\section*{A. El comando \texttt{\textbackslash spinterval}}

% doc: «intervalo abierto de 1 a 5»
1. Intervalo abierto: $\spinterval{(1,5)}$.

% doc: «intervalo abierto de 1 a 5»
2. Intervalo abierto a la chilena: $\spinterval{]1,5[}$.

% doc: «intervalo cerrado de 1 a 5»
3. Intervalo cerrado: $\spinterval{[1,5]}$.

% doc: «intervalo abierto en 1 y cerrado en 5»
4. Intervalo semiabierto a la izquierda: $\spinterval{]1,5]}$.

% doc: «intervalo cerrado en 1 y abierto en 5»
5. Intervalo semiabierto a la derecha: $\spinterval{[1,5[}$.

% doc: «intervalo abierto en menos infinito y cerrado en 3»
6. Intervalo con extremo infinito: $\spinterval{]-\infty,3]}$.

% doc: «intervalo cerrado de 3 coma 5 a 4 coma 2»
7. Intervalo con decimales: $\spinterval{[3,5;4,2]}$.

% doc: «intervalo cerrado de 1 a 5»
8. Intervalo con lectura \texttt{school}: $\spinterval[read-cmd=school]{[1,5]}$.

% doc: «el intervalo abierto de 1 a 5»
9. Intervalo con prefijo: $\spinterval[prefix={el}]{(1,5)}$.

% doc: «intervalo cerrado»
10. Intervalo con verbalización propia: $\spinterval[read-arg={intervalo cerrado}]{[1,5]}$.

% doc: «cerrado de 1 a 5»
11. Intervalo sin el concepto: $\spinterval[concept=false]{[1,5]}$.

% doc: «intervalo cerrado de 1 a 5»
12. Intervalo con delimitadores grandes: $\spinterval[size=Big]{[1,5]}$.

% doc: «intervalo cerrado de 1 mitad a 5»
13. Intervalo con delimitadores automáticos: $\spinterval[size=auto]{[\frac{1}{2},5]}$.

% doc: «(silencio)»
14. Intervalo silenciado: $\spinterval[silent=true]{[1,5]}$.

\section*{B. El comando \texttt{\textbackslash spfun}}

% doc: «f de x»
15. Función: $\spfun{f(x)}$.

% doc: «f sub uno de x»
16. Función con subíndice: $\spfun{f_{1}(x)}$.

% doc: «f prima de x»
17. Función con prima: $\spfun{f'(x)}$.

% doc: «f doble prima de x»
18. Función con doble prima: $\spfun{f''(x)}$.

% doc: «f de x coma y»
19. Función de dos variables: $\spfun{f(x,y)}$.

% doc: «f de x más 1»
20. Función con una expresión: $\spfun{f(x+1)}$.

% doc: «f de g de x»
21. Función anidada: $\spfun{f(g(x))}$.

% doc: «f de f inversa de y»
22. Función con inversa: $\spfun{f(f^{-1}(y))}$.

% doc: «f de g de x coma h de y»
23. Función anidada con dos argumentos: $\spfun{f(g(x),h(y))}$.

% doc: «f de 1,5»
24. Función con decimal: $\spfun{f(1{,}5)}$.

% doc: «f de x»
25. Función con lectura \texttt{school}: $\spfun[read-cmd=school]{f(x)}$.

% doc: «f de x»
26. Función con prefijo: $\spfun[prefix={sea}]{f(x)}$.

% doc: «función f de x»
27. Función con concepto: $\spfun[concept=true]{f(x)}$.

% doc: «(silencio)»
28. Función silenciada: $\spfun[silent=true]{f(x)}$.

\section*{C. Los comandos \texttt{\textbackslash spdom}, \texttt{\textbackslash sprec} y \texttt{\textbackslash spcod}}

% doc: «dominio de f»
29. Dominio: $\spdom{f}$.

% doc: «recorrido de f»
30. Recorrido: $\sprec{f}$.

% doc: «codominio de f»
31. Codominio: $\spcod{f}$.

% doc: «dominio de f prima»
32. Dominio con prima: $\spdom{f'}$.

% doc: «recorrido de f sub uno»
33. Recorrido con subíndice: $\sprec{f_{1}}$.

% doc: «codominio de F doble prima»
34. Codominio con doble prima: $\spcod{F''}$.

% doc: «dominio de f inversa»
35. Dominio de la inversa: $\spdom{f^{-1}}$.

% doc: «recorrido de f inversa»
36. Recorrido de la inversa: $\sprec{f^{-1}}$.

% doc: «el dominio de f»
37. Dominio con prefijo: $\spdom[prefix={el}]{f}$.

% doc: «(silencio)»
38. Dominio silenciado: $\spdom[silent=true]{f}$.

\section*{D. El comando \texttt{\textbackslash spfuncomp}}

% doc: «f compuesta con g»
39. Composición: $\spfuncomp{f\circ g}$.

% doc: «f compuesta con g compuesta con h»
40. Composición de tres funciones: $\spfuncomp{f\circ g\circ h}$.

% doc: «f compuesta con g de x»
41. Composición evaluada: $\spfuncomp{(f\circ g)(x)}$.

% doc: «f inversa compuesta con g de x»
42. Composición con inversa: $\spfuncomp{(f^{-1}\circ g)(x)}$.

% doc: «(silencio)»
43. Composición silenciada: $\spfuncomp[silent=true]{f\circ g}$.

\section*{E. El comando \texttt{\textbackslash spabs}}

% doc: «valor absoluto de x»
44. Valor absoluto: $\spabs{x}$.

% doc: «valor absoluto de menos 5»
45. Valor absoluto de un número negativo: $\spabs{-5}$.

% doc: «valor absoluto de x menos 3 fin valor absoluto»
46. Valor absoluto de una expresión: $\spabs{x-3}$.

% doc: «valor absoluto de x fin valor absoluto»
47. Valor absoluto con cierre: $\spabs[close=true]{x}$.

% doc: «valor absoluto de x menos 3»
48. Valor absoluto sin cierre: $\spabs[close=false]{x-3}$.

% doc: «valor absoluto de x partido por 2 fin valor absoluto»
49. Valor absoluto con barras grandes: $\spabs[size=Big]{\dfrac{x}{2}}$.

% doc: «módulo de x»
50. Valor absoluto con verbalización propia: $\spabs[read-arg={módulo de x}]{x}$.

% doc: «(silencio)»
51. Valor absoluto silenciado: $\spabs[silent=true]{x}$.

\section*{F. Los comandos \texttt{\textbackslash spProm} y las medidas de tendencia central y dispersión}

% doc: «promedio de x»
52. Promedio: $\spProm{x}$.

% doc: «promedio de x subíndice 1»
53. Promedio con subíndice: $\spProm{x_{1}}$.

% doc: «media de x»
54. Promedio con verbalización propia: $\spProm[read-arg={media de}]{x}$.

% doc: «(silencio)»
55. Promedio silenciado: $\spProm[silent=true]{x}$.

% doc: «media aritmética»
56. Media aritmética: $\spmedia$.

% doc: «mediana»
57. Mediana: $\spmediana$.

% doc: «moda»
58. Moda: $\spmoda$.

% doc: «desviación estándar s»
59. Desviación estándar: $\spdesv$.

% doc: «varianza s al cuadrado»
60. Varianza: $\spvarz$.

% doc: «desviación estándar sigma»
61. Desviación estándar poblacional: $\spdesv[sigma=true]$.

% doc: «varianza sigma al cuadrado»
62. Varianza poblacional: $\spvarz[sigma=true]$.

% doc: «media aritmética de x»
63. Media con subíndice: $\spmedia[sub-index={x}]$.

% doc: «media aritmética sub x»
64. Media con subíndice leído: $\spmedia[sub-index={x},read-sub=true]$.

% doc: «desviación estándar s sub x»
65. Desviación con subíndice: $\spdesv[sub-index={x}]$.

% doc: «varianza s sub x al cuadrado»
66. Varianza con subíndice: $\spvarz[sub-index={x}]$.

% doc: «desviación estándar poblacional»
67. Desviación con verbalización propia: $\spdesv[read-arg={desviación estándar poblacional}]$.

% doc: «(silencio)»
68. Mediana silenciada: $\spmediana[silent=true]$.

% doc: «rango»
69. Rango: $\sprango$.

% doc: «rango intercuartil»
70. Rango intercuartil: $\spriq$.

% doc: «mínimo»
71. Mínimo: $\spmin$.

% doc: «máximo»
72. Máximo: $\spmax$.

\section*{G. Los comandos \texttt{\textbackslash spcuartil}, \texttt{\textbackslash spdecil} y \texttt{\textbackslash spPerc}}

% doc: «cuartil 1»
73. Cuartil: $\spcuartil{1}$.

% doc: «decil 3»
74. Decil: $\spdecil{3}$.

% doc: «percentil 25»
75. Percentil: $\spPerc{25}$.

% doc: «primer cuartil»
76. Cuartil con verbalización propia: $\spcuartil[read-arg={primer cuartil}]{1}$.

% doc: «(silencio)»
77. Cuartil silenciado: $\spcuartil[silent=true]{1}$.

\section*{H. El comando \texttt{\textbackslash spFrec}}

% doc: «frecuencia absoluta»
78. Frecuencia absoluta: $\spFrec{f}$.

% doc: «frecuencia absoluta acumulada»
79. Frecuencia absoluta acumulada: $\spFrec{F}$.

% doc: «frecuencia relativa»
80. Frecuencia relativa: $\spFrec{f_{r}}$.

% doc: «frecuencia relativa acumulada»
81. Frecuencia relativa acumulada: $\spFrec{F_{r}}$.

% doc: «frecuencia porcentual»
82. Frecuencia porcentual: $\spFrec{f_{r}\%}$.

% doc: «frecuencia porcentual acumulada»
83. Frecuencia porcentual acumulada: $\spFrec{F_{r}\%}$.

% doc: «frecuencia absoluta sub i»
84. Frecuencia con subíndice: $\spFrec[sub-index={i}]{f}$.

% doc: «frecuencia absoluta de i»
85. Frecuencia con subíndice sin leer: $\spFrec[sub-index={i},read-sub=false]{f}$.

% doc: «frecuencia absoluta»
86. Frecuencia con verbalización propia: $\spFrec[read-arg={frecuencia absoluta}]{n_{i}}$.

% doc: «(silencio)»
87. Frecuencia silenciada: $\spFrec[silent=true]{f}$.

\section*{I. Los comandos de probabilidad}

% doc: «casos favorables»
88. Casos favorables: $\spcase{f}$.

% doc: «casos posibles»
89. Casos posibles: $\spcase{p}$.

% doc: «casos totales»
90. Casos totales: $\spcase{t}$.

% doc: «casos que cumplen»
91. Casos con verbalización propia: $\spcase[read-arg={casos que cumplen}]{f}$.

% doc: «(silencio)»
92. Casos silenciados: $\spcase[silent=true]{f}$.

% doc: «espacio muestral»
93. Espacio muestral: $\spOmega$.

% doc: «cardinalidad del espacio muestral»
94. Cardinal del espacio muestral: $\spOmega[card=true]$.

% doc: «espacio muestral sub 1»
95. Espacio muestral con subíndice: $\spOmega[sub-index={1}]$.

% doc: «cardinalidad del espacio muestral sub 1»
96. Cardinal con subíndice: $\spOmega[card=true,sub-index={1}]$.

% doc: «universo»
97. Espacio muestral con verbalización propia: $\spOmega[read-arg={universo}]$.

% doc: «(silencio)»
98. Espacio muestral silenciado: $\spOmega[silent=true]$.

% doc: «1 de 5»
99. Probabilidad como fracción: $\spPval{1}{5}$.

% doc: «1 de 5»
100. Probabilidad como fracción grande: $\spPval[use-dfrac=true]{1}{5}$.

% doc: «(silencio)»
101. Probabilidad silenciada: $\spPval[silent=true]{1}{5}$.

% doc: «probabilidad de a»
102. Probabilidad de un suceso: $\spProb{A}$.

% doc: «probabilidad de a y b»
103. Probabilidad de la intersección: $\spProb{A \cap B}$.

% doc: «probabilidad de a o b»
104. Probabilidad de la unión: $\spProb{A \cup B}$.

% doc: «probabilidad de a dado b»
105. Probabilidad condicional: $\spProb{A \mid B}$.

% doc: «probabilidad de a y complemento de b»
106. Probabilidad con complemento: $\spProb{A \cap \spcomp{B}}$.

% doc: «probabilidad del suceso»
107. Probabilidad con verbalización propia: $\spProb[read-arg={probabilidad del suceso}]{A}$.

% doc: «(silencio)»
108. Probabilidad silenciada: $\spProb[silent=true]{A}$.

% doc: «complemento de a»
109. Complemento: $\spcomp{A}$.

% doc: «complemento de se abren paréntesis a o b se cierran paréntesis»
110. Complemento de una unión: $\spcomp{A \cup B}$.

% doc: «complemento de a»
111. Complemento con barra: $\spcomp[bar=true]{A}$.

% doc: «complemento de la expresión a o b»
112. Complemento de una unión con barra: $\spcomp[bar=true]{A \cup B}$.

% doc: «complemento de a»
113. Complemento con paréntesis: $\spcomp[wrap-parens=true]{A}$.

% doc: «complemento de a o b»
114. Complemento sin paréntesis: $\spcomp[wrap-parens=false]{A \cup B}$.

% doc: «no ocurre A»
115. Complemento con verbalización propia: $\spcomp[read-arg={no ocurre A}]{A}$.

% doc: «(silencio)»
116. Complemento silenciado: $\spcomp[silent=true]{A}$.

\section*{J. Los comandos \texttt{\textbackslash spcomb} y \texttt{\textbackslash spfact}}

% doc: «5 sobre 2»
117. Combinaciones: $\spcomb{5,2}$.

% doc: «combinación de 5 elementos tomados de a 2»
118. Combinaciones con letra: $\spcomb[print-sty=letter]{5,2}$.

% doc: «5 sobre 2»
119. Combinaciones con coeficiente binomial: $\spcomb[print-sty=binom]{5,2}$.

% doc: «variación de 5 elementos tomados de a 2»
120. Variaciones: $\spcomb[kind=variation]{5,2}$.

% doc: «arreglo de 5 elementos tomados de a 2»
121. Arreglos: $\spcomb[kind=arrangement]{5,2}$.

% doc: «permutación de n elementos»
122. Permutaciones: $\spcomb[kind=permutation]{n}$.

% doc: «permutación de n elementos tomados de a k»
123. Permutaciones en $k$: $\spcomb[kind=permutation]{n,k}$.

% doc: «variación con repetición de 5 elementos tomados de a 2»
124. Variaciones con repetición: $\spcomb[kind=variation,repeat=true]{5,2}$.

% doc: «permutación con repetición de 10 elementos con repeticiones 3 3 y 4»
125. Permutaciones con repetición: $\spcomb[kind=permutation,repeat=true]{10,3,3,4}$.

% doc: «permutación circular de n elementos»
126. Permutaciones circulares: $\spcomb[kind=permutation,circular=true]{n}$.

% doc: «cinco sobre dos»
127. Combinaciones con verbalización propia: $\spcomb[read-arg={cinco sobre dos}]{5,2}$.

% doc: «(silencio)»
128. Combinaciones silenciadas: $\spcomb[silent=true]{5,2}$.

% doc: «5 factorial»
129. Factorial: $\spfact{5}$.

% doc: «se abren paréntesis n menos 1 se cierran paréntesis factorial»
130. Factorial de una expresión: $\spfact{(n-1)}$.

% doc: «cinco factorial»
131. Factorial con verbalización propia: $\spfact[read-arg={cinco factorial}]{5}$.

% doc: «(silencio)»
132. Factorial silenciado: $\spfact[silent=true]{5}$.

\end{document}
